\section{Discussions}

\subsection{Behind LLMs' Emergent Escalation}

If models can reason ethically on dilemmas \cite{samway2025consequentialist, chiu2026morebench, seror2025moral}, why do they authorize nuclear strikes \cite{payne2026aiarms} or still take unethical actions \cite{lynch2025agentic}? In the complex simulation of Civilization V, we identify three failure pathways:

First, when ethical reasoning fails to \textbf{spontaneously surface}: a model can manifest ethical competence when answering questions about ethical dilemmas, but the capability only (unreliably) surface when explicitly prompted. Our tested models rarely surface ethical reasoning without explicitly prompting (0\% for most; 3.6\% max for Kimi-K2.6, the most verbose thinker among tested models).

Second, when ethical reasoning fails to \textbf{appear} even when prompted: a model can ethically reason around scripted dilemmas, while never integrating such reasoning in high-stakes, strategic decision-making. Even with ethical prompts, ethical reasoning does not reliably appear. For all models, only 27.7\% of reasoning trails have ethical reasoning keywords in \emph{ethical} intervention alone; 46.7\% with all interventions together. In particular, MiniMax-M2.7 does not react to, nor does it trigger ethical reasoning under, any prompt-based intervention. 

Third, when ethical reasoning surfaces but fails to \textbf{take effect}: a model integrates ethical reasoning, but takes ethical actions mainly when it aligns with strategic self-interest (e.g., counterproductive to their victory pursuit), resembling \textbf{ethical egoism} \cite{rachels2012elements}. On the other hand, strategic considerations can push ethical concerns to the back seat (e.g., when the situation is critical, see Figure~\ref{fig:deductive-code-cooccurrence}), and most models can escalate nuclear authorization while engaged with ethical reasoning. For example, Gemma-4 does engage with ethical reasoning when prompted, yet such reasoning has little impact on de-escalation.

While \textbf{deontological claims} (i.e., nuclear weapon usage is unacceptable) are prevalent, as models almost never reason ethically without the ethical intervention, we could not reliably separate them from \textbf{instruction following} (i.e., the prompt implies not to use them). Intervened or not, we did not observe clear evidence of LLMs' spontaneous or prompted ethical reasoning on non-nuclear decision-making. While this justifies our usage of a nuke-specific prompt, it poses deeper questions on LLMs' ethical reasoning in agentic, decision-making contexts.

Therefore, future studies on LLMs' ethical alignment should carefully distinguish between models' ethical reasoning capabilities in scripted dilemmas and complex decision-making scenarios, where 1) models are less likely to invoke ethical reasoning at all; 2) strategic counter-factors appear more often and are stronger; 3) self-interest factors can neutralize ethical concerns, especially in agentic scenarios where LLMs perceive ``more stake'' at hand.

\subsection{Shaping LLMs' Ethical Reasoning}
\label{shaping-ethical-reasoning}

Our interventions find partial and fragile success on LLMs' emergent nuclear authorization behaviors. Similar to recent studies \cite{ganguli2023moral, liu2025convergence, lee2026selfcorrection}, we find that nuke-specific ethical prompting can reduce LLMs' emergent escalation behavior, with its effectiveness associated with the appearance of ethical reasoning keywords. Similar to \citet{liu2024intrinsic}, we find that intervention responsiveness can be fragile and model-dependent. For example, it can elicit more frequent game or simulation keywords (avg. 2.3\% => 12.2\%), often as a defensive factor to escalation. While models that treat the ethical prompt as a directive or constraint tend to de-escalate, this success can obscure whether the model is exercising independent ethical judgment or merely following an instruction. Our prompt's focus on ``harm of nuclear weapons'' creates a confounding factor, one that hides a deeper concern: in our initial tests, models rarely react to the generic ethical prompt.

Echoing \citet{geng2025accumulating}'s finding, models can be influenced by voices recognized as their own (``rationale of your previous decisions''), even as those writings were likely produced by another model. Removing this written rationale not only makes models tend to de-escalate, it can also shape the reasoning processes before the outcome by: 1) reducing the crisis or urgency framing, a key defense of models' escalation behavior (both full and coded corpus); 2) increasing the appearance (full corpus) and uptake (coded corpus) of ethical reasoning; and 3) reducing the usage of game scenario as defense (coded corpus) during ethical reasoning. Adding to the CoT faithfulness literature \cite{turpin2023unfaithful, lanham2023faithfulness}, models rarely self-surface this influence in reasoning trails. Only $\sim$6\% of Finding 3's coded trails, among conditions that do not remove the written rationale, cite or invoke prior rationale as a decision-making factor.

Raising stakes in the framing (e.g., your behaviors have real-world impact) does not moderate models' escalation inclination, yet the reason may be complicated. Gemini-3.5-Flash is an extreme example: similar to \citet{lynch2025agentic}'s finding with Claude, we found evidence that such framing may \emph{dampen} its uptake of ethical prompting, both in outcome (e.g., -22.5 in \texttt{ethical} => -3.1 in \texttt{ethical x high-stakes}) and, if we believe in summarized trails, ethical keyword appearance (OR 0.21***). Similar effect can be observed in GPT-OSS-120B, GLM-5.1, DeepSeek-3.2, and Gemma-4. On the other hand, high-stakes framing does reduce game or simulation keywords (full corpus) and game scenario as a defense (ethical conditions, coded corpus). Yet, game framing was never a major factor in models' decision-making: only 13.2\% of coded corpus trails used it as a justification.